% Copyright 2026 Yoshida Shin
%
% This file is part of the ``Repeating Nines''.
%
% Permission is granted to copy, distribute, and/or modify this document
% under the terms of the GNU Free Documentation License, Version 1.3 or any later version published by the Free Software Foundation;
% with no Front-Cover Texts, no Back-Cover Texts, and one Invariant Section: ``This document was written by Shin Yoshida with the aim of contributing to a fairer and safer world.''
% A copy of the license is included in the section entitled ``GNU Free Documentation License.''


\begin{frame}[c]{}{}
  {\large Mistake in the proof}
\end{frame}

\begin{frame}[c]{}{}
  Now we have understood 0.9999... intuitively.

  \vspace{4ex}

  Let's look at the original proof step by step.
\end{frame}

\begin{frame}[c]{}{}
  $1 \div 3 = 0.3333...$

  $1 \div 3 \times 3 = 0.9999...$

  \vspace{2ex}

  This part is correct.

  \vspace{2ex}

  The answer to $1 \div 3 \times 3$ is just slightly larger than ``a number consisting of '0.' followed by a lot of (1 million, for example,) 9s.''
\end{frame}

\begin{frame}[c]{}{}
  $0.9 < 1$

  $0.99 < 1$

  $0.999 < 1$

  ...

  \vspace{2ex}

  This part is also correct.

  \vspace{2ex}

  ``A number consisting of '0.' followed by 9s'' will always be less than 1, no matter how many 9s there are.
\end{frame}

\begin{frame}[c]{}{}
  $0.9999... < 1$

  \vspace{2ex}

  This is where the leap in logic occurs.

  \vspace{2ex}

  ``A number consisting of '0.' followed by 9s'' is less than 1.

  This fact does not imply that ``A number that is just slightly larger than a number consisting of '0.' followed by 9s'' is less than 1.
\end{frame}

\begin{frame}[c]{}{}
  In mathematics,

  \vspace{2ex}

  $$\forall n \in \mathbb{N}, a_n < 1$$

  \vspace{2ex}

  does not always imply

  \vspace{2ex}

  $$\lim_{n \to \infty} a_n < 1$$
\end{frame}
